\section{Finite balancing and Hermitian--Einstein metrics}\label{sec:metrics}

\subsection{The finite correction}
We recall the finite-dimensional result used below. Its independence
from weak orbit directions is the reason that the ambient estimate
of Proposition~\ref{prop:quadratic-form} suffices.

\begin{lemma}[{\cite[Theorem~2.1]{wuebben14}}]\label{lem:balancing}
Let $V$ be a finite quiver representation with compact group
$K=\prod_i\mathrm U(W_i)$ and flat moment map $\mu$ as in
\eqref{eq:moment}. On a ball suppose $\Omega$ is closed and invariant,
$\rho(0)=0$,
\[
 d\langle\rho,S\rangle=2\Omega(iC_xS,\cdot),\qquad
 \|\Omega-\omega_{\mathrm{flat}}\|\le\varepsilon .
\]
For sufficiently small $\varepsilon$, depending only on the fixed
multiplicity dimensions, a point with $\mu(x)=\zeta$ has a correction
$y=e^Sx$ with $\rho(y)=\zeta$ and $\|S\|_{\mathrm{op}}\le A\varepsilon$.
Conversely, a point with $\rho(y)=\zeta$ has a correction of the same
size to the flat level. Here $\zeta$ is central, and the starting
point lies in a smaller ball so that the correction stays in the domain.
Both assertions restrict to any locally $G$-invariant analytic subset.
They identify the two local compact quotients continuously, including
changes of stabilizer.
\end{lemma}

\begin{proof}
We recall the projection estimate and degree argument.
For an orthogonal vertex projection tuple $P$, write $x=n+u$, with
$n$ preserving and $u$ crossing the projected subspaces. On $u=0$
the infinitesimal field of $P$ vanishes, so both normalized moment
components paired with $P$ vanish. Integration along $n+su$ gives
\begin{equation}\label{eq:projection-error}
 |\langle\rho(x)-\mu(x),P\rangle|
                      \le\varepsilon\|C_xP\|^2 .
\end{equation}
The integration is in the ambient representation; its intermediate
points need not satisfy the obstruction equations.

Put $H=K_x$. Work in the Hermitian centralizer of $H$, perpendicular
to its Hermitian stabilizer algebra:
\[
 \p_H=\Herm\cap\operatorname{Cent}(H)\cap\mathfrak h^\perp.
\]
The pairing here is ordinary multiplicity trace.
On the connected $H$-fixed ball, the normalized moments annihilate
$\mathfrak h$ and are invariant under $H$, hence lie in $\p_H$.
The central target belongs to this space if either moment attains it
at the starting point.

Let $M=\sum_i m_i$. If $S\in\p_H$ has operator norm $a>0$, a
spectral projection above a global eigenvalue gap of size at least
$a/(M-1)$ satisfies $C_xP\ne0$. Otherwise $P$ and $I-P$ would fix
$x$ and be in $\mathfrak h$, while their trace pairings with $S$
vanish, contradicting the strict spectral gap. The identity is in
$\mathfrak h$, so $\tr S=0$. If $\p_H=0$ no correction is needed;
this also handles the trivial $M=1$ action.

Diagonalize $S$. A crossing coefficient of squared norm $w$ and
eigenvalue difference $d$ changes its squared norm by $e^{2d}w$.
Consequently
\begin{equation}\label{eq:spectral-cut}
 \begin{split}
 \langle\mu(e^Sx)-\mu(x),P\rangle
       &\ge 2e^{-4a}\frac{a}{M-1}\|C_xP\|^2,\\
 |\langle(\rho-\mu)(e^Sx),P\rangle|
       &\le\varepsilon e^{4a}\|C_xP\|^2 .
 \end{split}
\end{equation}
For $a=2M\varepsilon$ and sufficiently small $\varepsilon$, the
first expression dominates. Thus the homotopy
\[
 \mu(e^Sx)-\zeta+s(\rho-\mu)(e^Sx),\qquad 0\le s\le1,
\]
has no zero on the boundary of that ball in $\p_H$. At $s=0$ its
only zero is $S=0$, with positive derivative, since
\[
 \frac{d}{d\tau}\langle\mu(e^{\tau T}x),T\rangle
                   =2\|C_{e^{\tau T}x}T\|^2 .
\]
Its degree is one. The derivative is inverted only at this fixed
point to define degree; no uniform lower eigenvalue bound is used.
This proves forward correction. Starting from $\rho(y)=\zeta$,
use the homotopy
$\mu(e^Sy)-\mu(y)-s(\rho-\mu)(y)$ and the same estimate for reverse
correction.

The corrections stay in any locally invariant analytic subset.
For uniqueness on compact quotients, polar decomposition reduces two
representatives in the same complex orbit to $y$ and $e^Ty$. The
squared norm along $e^{sT}y$ is convex, so the path stays in the
ambient ball. Positivity gives
\[
 \frac{d}{ds}\langle\rho(e^{sT}y),T\rangle
        \ge2(1-\varepsilon)\|C_{e^{sT}y}T\|^2 .
\]
The endpoint moment values agree by centrality; hence $T$ fixes $y$.
The same argument works for $\mu$. Bounded correcting matrices,
compactness of $K$, and uniqueness of the resulting classes give
continuous inverse maps on nested local quotients.
\end{proof}

\subsection{Reconstruction of the metric}
The complementary equation makes $\M_t(\Psi_t(y))$ parallel.
By \eqref{eq:central-value}, its vanishing is equivalent to
\[
 \rho_t(y)=\zeta(t).
\]
Proposition~\ref{prop:quadratic-form} and Lemma~\ref{lem:balancing}
apply on the full representation $V$, then restrict to $Z$.
For $x\ne0$, use an ambient ball whose radius is a fixed multiple
of $\|x\|$. Its form error is at most $C(|t|+\|x\|^2)$.
We obtain
\begin{equation}\label{eq:S-bound}
 y=e^Sx\in Z,\qquad \rho_t(y)=\zeta(t),\qquad
             \|S\|_{\mathrm{op}}\le C\delta.
\end{equation}
At $x=0$, flat balance forces $\zeta=0$; the normalized corrected
moment also vanishes there, so take $S=0$.

Use the trace-free parallel lift $S^0$ in \eqref{eq:tracefree-lift}.
It acts on $x$ as $S$ does, preserves the determinant, and has norm
at most a fixed multiple of $\|S\|$. The actual metric is
\begin{equation}\label{eq:actual-metric}
 h^{\mathrm{HE}}_{t,x}
    =b_t(e^{S^0})^*\bigl(h_0e^{2\sigma(t,y)}\bigr).
\end{equation}
The star is pullback of a metric by the holomorphic isomorphism
$e^{S^0}:E_x\to E_y$.
The complementary estimate and $\|y\|\le e^{2\|S\|}\|x\|$
give \eqref{eq:metric-main}. Multiplication by $b_t$ supplies the
scalar determinant equation, hence the full HE equation.

\subsection{Coverage of nearby HYM solutions}\label{subsec:coverage}
We specify the qualitative inputs and their compatibility with the
quantitative construction. In the integrable HYM specialization of
\cite{delloque}, Section~3.2 makes the subsolution condition automatic.
Proposition~4.3 gives the convergent harmonic slice and local
classification; Proposition~4.5 gives the complementary equation.
Corollary~5.2 identifies nearby holomorphic morphisms with parallel
intertwiners. Proposition~5.6 supplies a common-parameter inverse,
Lemma~7.3 its open integrable image, and Theorem~7.4 the local
quotient topology. These results permit the complete multiplicity
algebra in \eqref{eq:split}.

Use a sufficiently small $W^{k,2}$ neighborhood of the central
operator, with $W^{k+1,2}$ gauges. The usual elliptic slice argument
gives the same maps in these Sobolev spaces; alternatively one may
apply the cited smooth or H\"older slice at higher regularity and
then use elliptic regularity. All actual reconstructed solutions
are smooth. On the fixed Hermitian bundle $h_0$, the representative
of a full solution is
\begin{equation}\label{eq:scalar-corrected}
 \widehat\Psi_t(y)=b_t^{1/2}\cdot\Psi_t(y),\qquad
       y\in Z,\quad \rho_t(y)=\zeta(t).
\end{equation}
Its determinant connection is the normalized HE determinant connection
at $\omega_t$, transported to $\det h_0$, and is independent of $y$.
We use determinant-preserving unitary gauges, up to the ineffective
constant scalar. Removing the known central gauge $b_t^{1/2}$
before taking the projective slice and restoring it afterward is
essential: $\det h_0$ itself need not be HE for $\omega_t$.
The scalar correction is smooth in $t$ and invisible on $\End E_0$.

For completeness, the passage from an arbitrary nearby positive
gauge to our complementary one uses no varying small eigenvalue.
If its logarithm is $s$, change the harmonic representative by a
parallel $e^h$ and consider
\begin{equation}\label{eq:positive-normalization}
 L(s,h)=\tfrac12\log(e^{-h}e^{2s}e^{-h}).
\end{equation}
The derivative in $h$ of its parallel projection at the origin is
$-I$. The finite implicit-function theorem therefore makes $L$
complementary on a common neighborhood. Uniqueness in
Lemma~\ref{lem:complement} identifies it with $\sigma(t,y)$ after
unitary gauge. The local slice inverse cited above then identifies
the actual solution with \eqref{eq:scalar-corrected}. Reverse
finite balancing supplies a flat-balanced $x$ in the same orbit.

Choose successive coefficient radii, for example
$R_0/16,R_0/8,R_0/4,R_0/2$, inside the common ambient $R_0$ ball.
After shrinking the parameter and actual neighborhoods, every
construction stays inside its domain. This proves coverage of an
open neighborhood of actual solutions, not only existence for
a prescribed family. Parallel-intertwiner classification identifies
isomorphic harmonic representatives with the same local $G$-orbit.
Uniqueness of the unitary HE class and the finite uniqueness argument
identify precisely the same equivalence classes. The common slice
inverse and finite quotient continuity prove the local homeomorphism
in Theorem~\ref{thm:main}.

\subsection{Smooth local pieces}\label{subsec:pieces}
Fix a flat-balanced point $x_0$ and put $H=K_{x_0}$.
Polar convexity gives $G_{x_0}=H^{\C}$.
The correction $S_0$ in Lemma~\ref{lem:balancing} lies in the
centralizer of $H$, so $y_0=e^{S_0}x_0$ has the same stabilizer.
On the $H$-fixed subspace the map
\begin{equation}\label{eq:piece-ift}
 (T,u)\longmapsto\rho_t(e^Te^{S_0}u)-\zeta(t),
                       \qquad T\in\p_H,
\end{equation}
takes values in $\p_H$. Its $T$ derivative at $(0,x_0)$ has positive
quadratic form
$2\Omega_t(iC_{y_0}T,C_{y_0}T)$, hence is invertible at that point.
The finite implicit-function theorem gives a smooth local correction.
Its domain may shrink with the least positive orbit eigenvalue;
the uniform estimate uses only the bounded value $S_0$, not the
derivative of this correction.

The ambient holomorphic slice theorem
\cite[Proposition~1.6 and Theorems~1.12, 1.20]{sjamaar}
applies to $V$ at a central moment value. Indeed the symplectic form
on the compact orbit is the pairing of $\zeta$ with Lie brackets,
so the orbit is isotropic. This works also at nonzero irrational
central values. Intersect the ambient slice with the reduced
analytic space $Z$. In the unchanged-stabilizer locus, the quotient
germ is represented by the $H$-fixed part of the transverse slice.
That analytic subset can still be singular. Refine it into smooth
analytic pieces and refine their boundaries by a local Whitney
stratification. These pieces have smooth lifts; the compact slice
and \eqref{eq:piece-ift} give balanced lifts.

Repeating at the other stabilizers yields a local cover by such
pieces, using the finite orbit types of the ambient compact linear
action and a locally finite analytic refinement.
Both quotient constructions use the same complex orbits and the
same pieces. This proves their smooth identification. It neither
identifies stabilizer type with smoothness nor asserts a uniform
lower radius for these individual charts.

